\documentclass[12pt,a4paper]{article}
\usepackage[utf8]{inputenc} % inutile avec XeLaTeX/LuaLaTeX
\usepackage[T1]{fontenc}
\usepackage{amsmath,amssymb,mathrsfs,tikz,times,pifont}
\usepackage{enumitem}
\usepackage{multicol}
\usepackage{lmodern}
\newcommand\circitem[1]{%
\tikz[baseline=(char.base)]{
\node[circle,draw=gray, fill=red!55,
minimum size=1.2em,inner sep=0] (char) {#1};}}
\newcommand\boxitem[1]{%
\tikz[baseline=(char.base)]{
\node[fill=cyan,
minimum size=1.2em,inner sep=0] (char) {#1};}}
\setlist[enumerate,1]{label=\protect\circitem{\arabic*}}
\setlist[enumerate,2]{label=\protect\boxitem{\alph*}}
\everymath{\displaystyle}
\usepackage[left=1cm,right=1cm,top=1cm,bottom=1.7cm]{geometry}
\usepackage[colorlinks=true, linkcolor=blue, urlcolor=blue, citecolor=blue]{hyperref}
\usepackage{array,multirow}
\usepackage[most]{tcolorbox}
\usepackage{varwidth}
\usepackage{float}
\tcbuselibrary{skins,hooks}
\usetikzlibrary{patterns}

\usepackage{pgfplots}
\pgfplotsset{compat=1.15}
\usepackage{mathrsfs}

\usepackage{amsmath}
\usepackage{amsfonts}
\usepackage{amssymb}
\usepackage{tkz-tab}
\usepackage{mdframed}
\usepackage{tikz}
\usetikzlibrary{arrows, shapes.geometric, fit}

% Commande pour la couleur d'accentuation
\newcommand{\myul}[2][black]{\setulcolor{#1}\ul{#2}\setulcolor{black}}
\newcommand\tab[1][1cm]{\hspace*{#1}}

%\usepackage[margin=2cm]{geometry}
\usepackage{eso-pic}         % Pour ajouter des éléments en arrière-plan

\usepackage{enumitem}
%---------------------------------------
% Définir un compteur pour les exemples
\newcounter{exemple}

% Définir la commande \exemple pour afficher un exemple numéroté
\newcommand{\exemple}{%
  \refstepcounter{exemple}% Incrémenter le compteur
  \textbf{\textcolor{orange}{Exemple \theexemple : }} \ignorespaces
}
%---------------------------------------
\newcounter{solution}

% Définir la commande \solutione pour afficher un solution numéroté
\newcommand{\solution}{%
  \refstepcounter{solution}% Incrémenter le compteur
  \textbf{\textcolor{orange}{Solution \thesolution : }} \ignorespaces
}
%---------------------------------------
\definecolor{myorange}{rgb}{1.0, 0.8, 0.0}

% Définir un compteur pour les exercices d'application
\newcounter{exerciceapp}

% Définir la commande pour afficher un exercice d'application numéroté
\newcommand{\exerciceapp}{%
  \refstepcounter{exerciceapp}%
  \textbf{\textcolor{myorange}{Exercice d'application \theexerciceapp :}} \ignorespaces
}
%--------------------------------------
% Définir un compteur pour les exercices d'application
\newcounter{correction}

% Définir la commande pour afficher un correction exercice d'application numéroté
\newcommand{\correction}{%
  \refstepcounter{correction}%
  \textbf{\textcolor{myorange}{Correction \thecorrection :}} \ignorespaces
}
%--------------------------------------
% Définir un compteur pour les remarque d'application
\newcounter{remarque}

%----------------------------------------
\definecolor{myorange1}{rgb}{1.0, 1.5, 0}
% Définir la commande pour afficher un remarque numéroté
\newcommand{\remarque}{%
  \refstepcounter{remarque}%
  \textbf{\textcolor{myorange1}{Remarque \theremarque :}} \ignorespaces
}

% Commande pour ajouter du texte en arrière-plan
\AddToShipoutPicture{
    \AtTextCenter{%
        \makebox[0pt]{\rotatebox{45}{\textcolor[gray]{0.6}{\fontsize{5cm}{5cm}\selectfont PGB}}}
    }
}

\begin{document}
% En-tête personnalisée
\begin{center}
    \Large\textbf{\underline{\textcolor{red}{Trinôme du $2^{\text{nd}}$ degré et système d’équations}}}\\[-0.1cm]
    \normalsize\textbf{Prof : M. BA} \hfill \textbf{Classe : TS2}\\[-0.1cm]
    \textbf{Année scolaire : 2025 -- 2026}
\end{center}

\section*{\underline{\textbf{\textcolor{red}{I.Trinôme du $2^{\text{nd}}$ degré}}}}

\subsection*{\underline{\textbf{\textcolor{red}{1. Définition}}}}
Soient $a,b,c \in \mathbb{R}$ tels que $a \neq 0$.

On appelle \textcolor{red}{trinôme du second degré} à variable $x$ toute expression qui s’écrit sous la forme :

\(
ax^2 + bx + c \qquad (a \neq 0)
\)
$a$, $b$ et $c$ sont appelés les coefficients du trinôme.

\textbf{\underline{\exemple}}

\begin{enumerate}
\item $H(x)=x^2-4$ est un trinôme du $2^{\text{nd}}$ degré tel que 

\(
a=1,\quad b=0 \quad \text{et} \quad c=-4.
\)

\item $G(x)=3x^2-x-1$ est un trinôme du $2^{\text{nd}}$ degré tel que
 
\(
a=3,\quad b=-1 \quad \text{et} \quad c=-1.
\)

\item $R(x)=x+4 \quad (R(x)\ \text{n’est pas un trinôme du } 2^{\text{nd}}\ \text{degré}).$

\end{enumerate}

\subsection*{\underline{\textbf{\textcolor{red}{2. Forme canonique d’un trinôme du $2^{\text{nd}}$ degré }}}}

\textcolor{red}{\underline{Activité}}

\begin{enumerate}
\item Soit $\alpha \in \mathbb{R}$. Montrer que pour tout réel $x$ on a :
\(
x^2+\alpha x=\left(x+\frac{\alpha}{2}\right)^2-\left(\frac{\alpha}{2}\right)^2
\)

\item Soient $P$, $G$ et $R$ des trinômes du $2^{\text{nd}}$ degré tels que :

\(
P(x)=2x^2-x+2
\)

\(
G(x)=3x^2-2x-1
\)

\(
R(x)=3x^2-2\sqrt{3}\,x+1
\)

En utilisant la relation démontrée dans 1), montrer que :

\begin{itemize}
\item[a)] $P(x)=2\left[\left(x-\frac{1}{2}\right)^2+\frac{15}{16}\right]$
\item[b)] $G(x)=3\left[\left(x-\frac{1}{3}\right)^2-\frac{4}{9}\right]$
\item[c)] $R(x)=3\left(x-\frac{\sqrt{3}}{3}\right)^2$
\end{itemize}

\end{enumerate}


\noindent
\textcolor{red}{\underline{\textbf{Théorème}}}

\medskip

Soit 
\(
P(x)=ax^2+bx+c
\)
où $a,b,c\in\mathbb{R}$ avec $a\neq0$.

Alors
\(
P(x)=a\left[\left(x+\frac{b}{2a}\right)^2-\frac{\Delta}{4a^2}\right]
\)
où
\(
\Delta=b^2-4ac.
\)

(On lit $\Delta$ « delta ».)

\medskip

L’expression obtenue dans $(*)$ est appelée \textcolor{red}{forme canonique} du trinôme du $2^{\text{nd}}$ degré $P(x)$.

\medskip

$\Delta$ est appelé le \textcolor{red}{discriminant} de $P(x)$.

\bigskip
\textcolor{red}{\underline{\textbf{Preuve}}}

\medskip

On part de l’expression :
\(
P(x)=ax^2+bx+c.
\)

On factorise par $a$ les termes contenant $x$ :
\(
P(x)=a\left(x^2+\frac{b}{a}x\right)+c.
\)

On complète le carré :
\(
x^2+\frac{b}{a}x
=
\left(x+\frac{b}{2a}\right)^2-\left(\frac{b}{2a}\right)^2.
\)

Donc
\(
P(x)
=
a\left[\left(x+\frac{b}{2a}\right)^2-\frac{b^2}{4a^2}\right]+c.
\)

On distribue $a$ :
\(
P(x)
=
a\left(x+\frac{b}{2a}\right)^2-\frac{b^2}{4a}+c.
\)

On met le second membre sous un seul dénominateur :
\(
-\frac{b^2}{4a}+c
=
\frac{-b^2+4ac}{4a}
=
-\frac{b^2-4ac}{4a}.
\)

Donc
\(
P(x)
=
a\left(x+\frac{b}{2a}\right)^2
-
\frac{b^2-4ac}{4a}.
\)

On factorise par $a$ :
\(
P(x)
=
a\left[
\left(x+\frac{b}{2a}\right)^2
-
\frac{b^2-4ac}{4a^2}
\right].
\)

Or
\(
\Delta=b^2-4ac.
\)

Ainsi
\(
P(x)=a\left[\left(x+\frac{b}{2a}\right)^2-\frac{\Delta}{4a^2}\right].
\)

\hfill $\square$

\textbf{\underline{\exerciceapp}}

Déterminer la forme canonique des trinômes du 2\textsuperscript{nd} degré suivants :

\begin{enumerate}
\item $P(x)=4x^2-2x+8$
\item $Q(x)=-3x^2+2\sqrt{6}\,x+2$
\item $R(x)=5x^2+4x-1$
\end{enumerate}

\textbf{\underline{\correction}}

\subsection*{\underline{\textbf{\textcolor{red}{3.Factorisation d’un trinôme du second degré}}}}
\subsection*{\underline{\textbf{\textcolor{red}{Définition}}}}

Soit $P$ un trinôme du second degré.  
Un réel $\alpha$ est \textcolor{red}{racine} de $P$ si $P(\alpha)=0$.

\textbf{\underline{\exemple}}

Soit 
\(
P(x)=3x^2-x-2.
\)

\begin{enumerate}
\item 
\(
P(1)=3(1)^2-1-2
\)

\(
P(1)=3-1-2
\)

\(
P(1)=0
\)

Donc \textcolor{red}{$1$ est racine de $P$.}

\item 
\(
P(0)=3(0)^2-0-2
\)

\(
P(0)=-2\neq0
\)

Donc \textcolor{red}{$0$ n’est pas racine de $P$.}
\end{enumerate}

\subsection*{\underline{\textbf{\textcolor{red}{Théorème}}}}

\medskip

Soit 
\(
P(x)=ax^2+bx+c
\)
où $a,b,c\in\mathbb{R}$ avec $a\neq0$.

\(
\Delta=b^2-4ac
\)

\begin{enumerate}
\item[(i)] Si $\Delta<0$ alors $P$ n’admet pas de racine.

\item[(ii)] Si $\Delta>0$ alors $P$ admet deux racines distinctes $x_1$ et $x_2$ telles que :

\(
x_1=\frac{-b-\sqrt{\Delta}}{2a}
\quad \text{et} \quad
x_2=\frac{-b+\sqrt{\Delta}}{2a}.
\)

\item[(iii)] Si $\Delta=0$ alors $P$ admet une racine double $x_0$ telle que :

\(
x_0=\frac{-b}{2a}.
\)
\end{enumerate}
\bigskip
\textcolor{red}{\underline{\textbf{Preuve}}}

\medskip

On sait que
\[
P(x)=a\left[\left(x+\frac{b}{2a}\right)^2-\frac{\Delta}{4a^2}\right].
\]

\[
P(x)=0
\iff
a\left[\left(x+\frac{b}{2a}\right)^2-\frac{\Delta}{4a^2}\right]=0
\]

\[
\iff
\left(x+\frac{b}{2a}\right)^2-\frac{\Delta}{4a^2}=0
\quad \text{car } a\neq0
\]

\[
\iff
\left(x+\frac{b}{2a}\right)^2=\frac{\Delta}{4a^2}
\tag{$*$}
\]

Ainsi, on distingue 3 cas suivant les valeurs de $\Delta$ :

\begin{enumerate}

\item[(i)] Si $\Delta<0$ alors la relation $(*)$ est impossible car le carré d’un réel n’est jamais négatif.

\medskip

\item[(ii)] Si $\Delta>0$ alors $(*)$ devient :

\[
x+\frac{b}{2a}
=
\sqrt{\frac{\Delta}{4a^2}}
\quad \text{ou} \quad
x+\frac{b}{2a}
=
-\sqrt{\frac{\Delta}{4a^2}}
\]

\[
x+\frac{b}{2a}
=
\frac{\sqrt{\Delta}}{2a}
\quad \text{ou} \quad
x+\frac{b}{2a}
=
-\frac{\sqrt{\Delta}}{2a}
\]

\[
x
=
-\frac{b}{2a}
+
\frac{\sqrt{\Delta}}{2a}
\quad \text{ou} \quad
x
=
-\frac{b}{2a}
-
\frac{\sqrt{\Delta}}{2a}
\]

\[
x_1=\frac{-b+\sqrt{\Delta}}{2a}
\quad \text{et} \quad
x_2=\frac{-b-\sqrt{\Delta}}{2a}
\]

\medskip

\item[(iii)] Si $\Delta=0$ alors $(*)$ devient :

\[
\left(x+\frac{b}{2a}\right)^2=0
\]

\[
x+\frac{b}{2a}=0
\]

\[
x=-\frac{b}{2a}
\]

\end{enumerate}

\hfill $\square$

\noindent
\textcolor{red}{\underline{Conséquence : Factorisation de $P(x)$}}

\medskip

Soit $P(x)=ax^2+bx+c$, $a,b,c\in\mathbb{R}$ avec $a\neq0$.

\begin{enumerate}
\item[(i)] Si $\Delta<0$ alors $P(x)$ n’a pas de factorisation dans $\mathbb{R}$.

\item[(ii)] Si $\Delta>0$ alors $P(x)$ est factorisable telle que :

\(
\boxed{P(x)=a(x-x_1)(x-x_2)}
\)

\item[(iii)] Si $\Delta=0$ alors $P(x)$ est factorisable telle que :

\(
\boxed{P(x)=a(x-x_0)^2}
\)
\end{enumerate}

\textcolor{red}{\underline{Application}}

Dans chacun des cas suivants, factoriser $P(x)$ si possible :

\begin{enumerate}
\item $P(x)=3x^2-x-2$
\item $P(x)=-x^2+x-3$
\item $P(x)=3x^2+2\sqrt{6}\,x+2$
\end{enumerate}

\noindent
\textcolor{red}{\underline{Remarque :}}

\medskip

\textbf{Calcul des racines avec $\Delta'$}

\medskip

Soit 
\(
P(x)=ax^2+bx+c
\)
où $a,b,c\in\mathbb{R}$ avec $a\neq0$.

Si 
\(
b'=\frac{b}{2}
\)
alors
\(
\Delta'=(b')^2-ac.
\)

\medskip

Tel que :

\begin{enumerate}

\item[(i)] Si $\Delta'<0$ alors $P(x)$ n’admet pas de racine dans $\mathbb{R}$.

\item[(ii)] Si $\Delta'>0$ alors $P(x)$ admet deux racines distinctes $x_1$ et $x_2$ telles que :

\(
x_1=\frac{-b'-\sqrt{\Delta'}}{a}
\quad \text{et} \quad
x_2=\frac{-b'+\sqrt{\Delta'}}{a}.
\)

\end{enumerate}

\begin{enumerate}
\item[(iii)] Si $\Delta'=0$ alors $P(x)$ admet une racine double $x_0$ telle que :
\[
x_0=\frac{-b'}{a}.
\]
\end{enumerate}

\medskip

\textbf{En effet :}

\[
\Delta=b^2-4ac
\]

Si $b'=\frac{b}{2}$ alors $b=2b'$.

Donc
\[
\Delta=(2b')^2-4ac
\]

\[
\Delta=4(b')^2-4ac
\]

\[
\Delta=4\big((b')^2-ac\big)
\]

Posons
\[
\Delta'=(b')^2-ac.
\]

Alors
\[
\Delta=4\Delta'.
\]

\medskip

\begin{enumerate}
\item[(i)] Si $\Delta'<0$ alors $\Delta<0$.

\item[(ii)] Si $\Delta'>0$ alors
\[
\sqrt{\Delta}=\sqrt{4\Delta'}=2\sqrt{\Delta'}.
\]

Ainsi,
\[
x_1=\frac{-b-\sqrt{\Delta}}{2a}
\quad \text{et} \quad
x_2=\frac{-b+\sqrt{\Delta}}{2a}.
\]

Or $b=2b'$ et $\sqrt{\Delta}=2\sqrt{\Delta'}$, donc

\[
x_1=\frac{-2b'-2\sqrt{\Delta'}}{2a}
\quad \text{et} \quad
x_2=\frac{-2b'+2\sqrt{\Delta'}}{2a}.
\]

En simplifiant :

\[
x_1=\frac{-b'-\sqrt{\Delta'}}{a}
\quad \text{et} \quad
x_2=\frac{-b'+\sqrt{\Delta'}}{a}.
\]

\item[(iii)] Si $\Delta'=0$ alors $\Delta=0$.

\[
x_0=\frac{-b}{2a}
=\frac{-2b'}{2a}
=\frac{-b'}{a}.
\]

\end{enumerate}

\hfill $\square$

\end{document}
